% !atlas-meta
% id: isha-upanishad
% title: Isha Upanishad
% author: badarayana
% authorName: Various
% language: English
% sourceLanguage: Sanskrit
% translator: F. Max Müller
% translationYear: 1879
% source: https://en.wikisource.org/wiki/Sacred_Books_of_the_East/Volume_1/V%C3%A2gasaneyi-samhit%C3%A2-upanishad
% license: Public Domain (translator F. Max Müller d. 1900)
% status: complete
% !end-atlas-meta
\documentclass[12pt]{article}
\usepackage[utf8]{inputenc}
\usepackage{ebgaramond}
\usepackage[margin=1.2in]{geometry}
\usepackage{parskip}
\newcommand{\work}[1]{\begin{center}\LARGE\scshape #1\end{center}}
\newcommand{\attrib}[2]{\begin{center}\large #1\\[2pt]\small\itshape trans.\ #2\end{center}\vspace{1em}}
\newcommand{\para}[1][]{\par\noindent\ifx&#1&\else\textsuperscript{\footnotesize #1}~\fi}
\begin{document}
% !body-start
\work{Isha-Upanishad (V\^agasaneyi-Samhit\^a-Upanishad)}
\attrib{Anonymous (Vedic sages)}{F. Max M\"uller (public domain)}

\para[1] All this, whatsoever moves on earth, is to be hidden in the Lord (the Self). When thou hast surrendered all this, then thou mayest enjoy. Do not covet the wealth of any man!

\para[2] Though a man may wish to live a hundred years, performing works, it will be thus with him; but not in any other way: work will thus not cling to a man.

\para[3] There are the worlds of the Asuras covered with blind darkness. Those who have destroyed their self (who perform works, without having arrived at a knowledge of the true Self), go after death to those worlds.

\para[4] That one (the Self), though never stirring, is swifter than thought. The Devas (senses) never reached it, it walked before them. Though standing still, it overtakes the others who are running. Mâtarisvan (the wind, the moving spirit) bestows powers on it.

\para[5] It stirs and it stirs not; it is far, and likewise near. It is inside of all this, and it is outside of all this.

\para[6] And he who beholds all beings in the Self, and the Self in all beings, he never turns away from it.

\para[7] When to a man who understands, the Self has become all things, what sorrow, what trouble can there be to him who once beheld that unity?

\para[8] He (the Self) encircled all, bright, incorporeal, scatheless, without muscles, pure, untouched by evil; a seer, wise, omnipresent, self-existent, he disposed all things rightly for eternal years.

\para[9] All who worship what is not real knowledge (good works), enter into blind darkness: those who delight in real knowledge, enter, as it were, into greater darkness.

\para[10] One thing, they say, is obtained from real knowledge; another, they say, from what is not knowledge. Thus we have heard from the wise who taught us this.

\para[11] He who knows at the same time both knowledge and not-knowledge, overcomes death through not-knowledge, and obtains immortality through knowledge.

\para[12] All who worship what is not the true cause, enter into blind darkness: those who delight in the true cause, enter, as it were, into greater darkness.

\para[13] One thing, they say, is obtained from (knowledge of) the cause; another, they say, from (knowledge of) what is not the cause. Thus we have heard from the wise who taught us this.

\para[14] He who knows at the same time both the cause and the destruction (the perishable body), overcomes death by destruction (the perishable body), and obtains immortality through (knowledge of) the true cause.

\para[15] The door of the True is covered with a golden disk. Open that, O Pûshan, that we may see the nature of the True.

\para[16] O Pûshan, only seer, Yama (judge), Sûrya (sun), son of Pragâpati, spread thy rays and gather them! The light which is thy fairest form, I see it. I am what He is (viz. the person in the sun).

\para[17] Breath to air, and to the immortal! Then this my body ends in ashes. Om! Mind, remember! Remember thy deeds! Mind, remember! Remember thy deeds

\para[18] Agni, lead us on to wealth (beatitude) by a good path, thou, O God, who knowest all things! Keep far from us crooked evil, and we shall offer thee the fullest praise!
% !body-end
\end{document}
